\documentstyle[%


righttag%


,11pt%


,amssymb%


,russian%


,izvru%


]{amsart}





\newcommand{\tl}[3]{\medskip\noindent{\bf#1}\ #2 \dotfill #3 \par} 


\pagestyle{empty}


\begin{document}


%\tableofcontents


\par


\vskip 2cm


\par


\bigskip





Номер посвящен исследованиям по теории дифференциальных


и функционально-диф\-фе\-рен\-ци\-аль\-ных уравнений. 


Изучаются краевые задачи теории


управления. Предыдущие номера: \No\,5, 1993; \No\,6, 1994;


\No\,10, 1995; \No\,11, 1996, \No\,6, 1997; \No\,2, 1999.





\bigskip


\bigskip


\par


\centerline{\Large СОДЕРЖАНИЕ}


\bigskip


\bigskip





\tl{Абдрахманов В.Г., Сапрыкин Е.Ф., Смолин Ю.Н.}


{Об устойчивости решения задачи\newline Коши для периодического


функционально-дифференциального уравнения с распределен-\newline ным


запаздыванием}{3}


\tl{Бравый Е.И.}


{О необходимых условиях


отрицательности функции Грина двухточечной\newline задачи}{12}


\tl{Власов В.В.} 


{O свойствах системы экспоненциальных решений 


дифференциально-разност-\newline ных уравнений в пространствах 


Соболева}{23}


\tl{Жуковский Е.С.} 


{О представлении оператора Грина абстрактного 


функционально-диффе-\newline ренциального уравнения}{30}


\tl{Иванов А.Г.}


{Об одном свойстве почти периодического


интеграла, зависящего от\newline параметра}{34}


\tl{Муминов Г.М.} 


{О двухточечной краевой задаче для линейных 


дифференциальных уравне-\newline ний второго порядка 


с коэффициентами, принимающими значения на заданном множестве}{44}


\tl{Пашуткин Д.В.}


{О приводимости нелинейных дифференциальных уравнений


к блочно-тре-\newline угольному виду}{50}


\tl{Перцев Н.В.}


{Двусторонние оценки решений интегродифференциального уравнения, \newline


описывающего процесс кроветворения}{58}


\tl{Солдатов М.А., Митрякова Т.М., Ильичев В.А.}


{Отыскание решений линейных диф-\newline ференциально-разностных уравнений 


в вырожденном случае}{63}


\tl{Цалюк З.Б.}


{Структура резольвенты системы уравнений восстановления


с разностным\newline ядром}{71}





\vfill


\noindent\raisebox{2pt}{\copyright} Казанский госудаpственный унивеpситет


\end{document}


